\documentclass[10pt,a4paper]{amsart}
\usepackage[margin=19mm]{geometry}
\usepackage{amsmath,amssymb,mathtools}
\usepackage[scr=rsfs,cal=euler]{mathalfa}
\usepackage{xfrac}
\usepackage[unicode,hidelinks,bookmarks=false]{hyperref}
\newcommand{\ie}{i.e.\ }
\newcommand{\cf}{cf.\ }
\newcommand{\resp}{respectively\ }
\newcommand{\Subsection}{\subsection}
\providecommand{\nobreakdash}{\nobreak\hbox{-}}
\setlength{\emergencystretch}{2em}
\multlinegap0pt
\usepackage{mathtools}
\usepackage{amssymb}
\usepackage{tensor}
\usepackage[shortlabels]{enumitem}
\usepackage{tikz}
\usepackage{float}
\usepackage{bm}
\usepackage{stackengine}
\newtheorem{theorem}{Theorem}
\newcommand{\e}{\mathcal{E}}
\newcommand{\h}{\mathcal{H}}
\newcommand{\m}{\mathcal{M}}
\newcommand{\mC}{\mathcal{C}}
\newcommand{\mo}{\mathcal{O}}
\newcommand{\p}{\mathbb{P}}
\newcommand{\ri}{\rightarrow}
\newcommand{\ve}{\mathcal{V}ec}
\title{\MakeUppercase{Cycles in the universal moduli stack of bundles of rank two over genus two curves}}
\author{Shubham Saha}
\date{Source: arXiv:2604.11958v1 (CC BY 4.0)}
\begin{document}
\maketitle

\noindent\emph{Source and licence.} \MakeUppercase{Cycles in the universal moduli stack of bundles of rank two over genus two curves}. Version arXiv:2604.11958v1 (CC BY 4.0), distributed by arXiv under the Creative Commons Attribution 4.0 International licence, http://creativecommons.org/licenses/by/4.0/, as recorded in the arXiv licence metadata for this exact version and shown on the abstract page https://arxiv.org/abs/2604.11958v1. Full licence notice: \texttt{licenses/arxiv-2604.11958-CC-BY-4.0.txt}.\par\smallskip

\noindent\textit{Editorial scope.} Counted unit: the article's theorem thm:cu (source0.tex lines 449--451). The article's complete original proof of it is delivered with it (source0.tex lines 452--460, one \texttt{proof} environment of 9 lines). No mathematics is written, repaired or re-derived: the statement and the proof are the article's own lines, in the article's own notation, and no retained line is substituted or re-flowed.  No further source-local context is needed: every cross-reference of every retained line resolves inside the two delivered spans.  Nothing was omitted from any retained span, so the package carries no omission record.  No retained line contains a citation, so the wrapper carries no bibliography and no bibliography entry.  No retained line contains a figure environment, an image-inclusion command or drawing code, so no image is delivered and the package declares no asset dependency.  Editorial formatting only, in addition: an AMS \texttt{amsart} preamble that restates the article's own declarations and macros used by the retained lines, namely the environments \texttt{theorem} on separate counters, together with the standard packages those lines need.  The retained environments print in this document as Theorem 1 in order of appearance, whereas the article prints the same items with the numbering of its own full document.  The complete native source of the article as distributed in the arXiv e-print for this exact version is bundled unchanged as \texttt{sources/198372/source0.tex}.
\begin{theorem}\label{thm:cu}
The Hecke correspondence along the Weierstrass section induces an isomorphism $c_w:\p_{3}\xrightarrow{\sim} \p_{2}$.
\end{theorem}

\begin{proof}
	The Poincar\'e bundle $\e_3\ri \mathcal{C}_{2,w}\times_{\mathcal{M}_{2,w}} \ve_w^\h(2,3,2)$ pulls back via $(1\times \phi_3)$ to the following composition of surjections:
	$$(1\times \phi_3)^*\e_3\twoheadrightarrow \phi_3^*(\e_3|_{w\times{\mathcal{M}_{2,w}} \ve_w^\h(2,3,2)})\twoheadrightarrow \mo_{w\times{\mathcal{M}_{2,w}}\p_3}(1).$$
Let $\e_2$ be the kernel of the composed surjection $(1\times \phi_3)^*\e_3\twoheadrightarrow\mo_{w\times_{\mathcal{M}_{2,w}}\p_3}(1)$. We show that $\e_2\ri \mathcal{C}_{2,w}\times_{\mathcal{M}_{2,w}} \p_3$ is a family of vector bundles with the claimed properties.\\
The sheaves $(1\times \phi_3)^*\e_3,\mo_{w\times_{\mathcal{M}_{2,w}} \p_3}(1)$ are flat over $\p_3$. Therefore, 
$$\e_2|_{C\times p}\simeq ker(\e_3|_{C_p\times \phi_3(p)}\twoheadrightarrow \mo_{w_p}) \forall p\in |\p_3|.$$
It follows that $deg(\e_2|_{C_p\times \phi_3(p)})=2 $ for all $p\in \p_3$. Hence we a have a morphism $c_w:\p_3\ri \p_2$.
To see that $c_w$ is an isomorphism, we construct the inverse map. The Hecke modification of $\e_2$ along the section $w$ produces a family of degree $1$ bundles on $\mC_{2,w}\times_{\m_{2,w}}\p_2$, twisting by $\mo(w)$ defines the inverse map $\p_2\ri \p_3$.
\end{proof}

\end{document}
